\section{Preliminaries}
\label{sec:preliminaries}

\subsection{Notation}
\label{sec:prelim:notation}

\begin{itemize}[nosep,leftmargin=2em]
  \item $\lambda$ --- security parameter (set to $192$ for SHAKE-192s,
        $256$ for SHAKE-256s).
  \item $n$ --- committee size, $1 \le n \le 256$.
  \item $t$ --- reconstruction threshold, $1 \le t \le n$.
  \item $\mathrm{GF}(p)$ --- finite field of order $p$; throughout
        this paper, $p = 257$ is the smallest prime greater than
        $255$.
  \item $\textsc{seed\_size}$ --- byte-length of the SLH-DSA scheme
        seed: $96$ bytes for SHAKE-192s and SHAKE-192f; $128$ bytes
        for SHAKE-256s. This is the secret input to FIPS 205
        $\textsc{DeriveKey}$.
  \item $P_i$ --- party $i$, $1 \le i \le n$. Each party is identified
        by a 32-byte $\mathsf{NodeID}$.
  \item $x_i \in [1, 256]$ --- party $P_i$'s Shamir evaluation point,
        distinct per party.
  \item $c_i \in \{0,1\}^{8 \cdot \textsc{seed\_size}}$ --- party
        $P_i$'s DKG contribution.
  \item $S \in \{0,1\}^{8 \cdot \textsc{seed\_size}}$ --- joint master
        SLH-DSA seed.
  \item $\mathit{share}_i \in [0, 257)^{\textsc{seed\_size}}$ ---
        party $P_i$'s Shamir share of the master byte-sum (per byte
        position, modulo $p = 257$).
  \item $(pk, sk)$ --- the FIPS 205 SLH-DSA public/private keypair
        derived from $S$ via FIPS 205 $\textsc{KeyGen}$.
  \item $\sigma$ --- a FIPS 205 SLH-DSA signature (the wire bytes).
  \item $\mu$ --- the message being signed; $\mathit{ctx}$ --- the
        FIPS 205 context string ($0$--$255$ bytes).
  \item $\mathit{sid} \in \{0,1\}^{128}$ --- a session identifier
        binding one threshold-sign attempt across the network.
  \item $\kappa \in \mathbb{N}$ --- session attempt counter (for
        replay isolation between retries).
  \item $\textsc{cSHAKE256}$ --- the cSHAKE256 customisable hash from
        NIST SP 800-185 \cite{sp800185}, with function-name string
        ``Magnetar'' and per-context customisation tags.
  \item $\Vert$ --- byte concatenation; $\oplus$ --- bitwise XOR.
\end{itemize}

\subsection{FIPS 205 SLH-DSA review}
\label{sec:prelim:slhdsa}

FIPS 205 \cite{fips205} (formerly SPHINCS+ \cite{sphincs}) is the NIST
stateless hash-based signature standard. The Magnetar reference
implementation targets the FIPS 205 SHAKE parameter sets:
\begin{table}[H]
\centering
\caption{FIPS 205 SHAKE parameter sets exposed by Magnetar v0.1. Values
from FIPS 205 \cite{fips205} \S10.1 Table~2.}
\label{tab:fips205-params}
\begin{tabular}{lrrrrrr}
\toprule
Parameter set & NIST PQ Cat. & $n$ (WOTS+) & $h$ & $d$ & $|pk|$ & $|\sigma|$ \\
\midrule
SLH-DSA-SHAKE-192s & 3 & 24 B & 63 & 7 & 48 B & 16\,224 B \\
SLH-DSA-SHAKE-192f & 3 & 24 B & 66 & 22 & 48 B & 35\,664 B \\
SLH-DSA-SHAKE-256s & 5 & 32 B & 64 & 8 & 64 B & 29\,792 B \\
\bottomrule
\end{tabular}
\end{table}

The FIPS 205 construction is built from the following layered
components, top to bottom:
\begin{itemize}[leftmargin=2em,nosep]
  \item \textbf{XMSS-HT (hypertree)}: a tree of $d$ XMSS layers. The
        leaves of the root layer are XMSS public keys; the leaves of
        the bottom layer are FORS public keys. Every layer has height
        $h/d$.
  \item \textbf{XMSS (eXtended Merkle Signature Scheme)}: a binary
        Merkle tree where each leaf is a WOTS+ public key.
  \item \textbf{WOTS+ (Winternitz One-Time Signature)}: a one-time
        signature whose security reduces to hash function preimage
        and second-preimage resistance.
  \item \textbf{FORS (Forest of Random Subsets)}: the few-time
        signature scheme used at the bottom of the hypertree.
  \item \textbf{Tweakable hash $\mathbf{F}, \mathbf{H}, \mathbf{T}_\ell$}:
        the cryptographic primitives, instantiated as SHAKE256 for the
        SHAKE-family Magnetar targets.
\end{itemize}

\paragraph{Key generation.} $\textsc{KeyGen}$ \cite{fips205}~\S9.1
samples a $3n$-byte secret triple $(\mathit{SK.seed},
\mathit{SK.prf}, \mathit{PK.seed})$ and computes
$\mathit{PK.root} = \mathbf{T}_h(\mathit{SK.seed}, \mathit{PK.seed},
\textsc{ADRS})$ as the XMSS-HT root. The public key is
$pk = (\mathit{PK.seed}, \mathit{PK.root})$; the secret key is
$sk = (\mathit{SK.seed}, \mathit{SK.prf}, \mathit{PK.seed},
\mathit{PK.root})$. The total scheme-seed size (the input to
$\textsc{DeriveKey}$) is $4n$ bytes: $96$ for $n=24$ (192s/192f) and
$128$ for $n=32$ (256s).

\paragraph{Deterministic signing.} $\textsc{Sign}_{\mathrm{det}}(sk, \mu,
\mathit{ctx})$ \cite{fips205}~\S9.2 computes a randomiser
$R \gets \mathrm{PRF}_{\mathit{SK.prf}}(\mathit{ctx} \Vert \mu)$ and
proceeds deterministically. The output is byte-identical across calls
with the same $(sk, \mu, \mathit{ctx})$. This determinism is the
property that makes the Magnetar Class-N1-analog claim trivially
provable: if two parties reconstruct the same $sk$ and call
$\textsc{Sign}_{\mathrm{det}}$ on the same $(\mu, \mathit{ctx})$, they
obtain byte-identical signatures.

\paragraph{Randomised signing.} $\textsc{Sign}_{\mathrm{rand}}(sk,
\mathit{rng}, \mu, \mathit{ctx})$ \cite{fips205}~\S9.2 samples a fresh
$n$-byte randomiser from $\mathit{rng}$. Magnetar's threshold-sign path
uses the deterministic variant by default (so the byte-equality claim
holds); a $\mathit{randomized}$ flag exposes the randomised variant
for callers that want hedge against fault-injection on the
$\mathit{PRF}$ derivation.

\paragraph{Verification.} $\textsc{Verify}(pk, \mu, \sigma,
\mathit{ctx})$ \cite{fips205}~\S9.3 walks the XMSS-HT root, recomputes
$\mathit{PK.root}$ from the FORS and XMSS authentication paths in
$\sigma$, and accepts iff it matches the $\mathit{PK.root}$ embedded
in $pk$. The verifier consumes no secret material; \emph{any
unmodified FIPS 205 verifier accepts a Magnetar signature with no
code change.}

\subsection{Shamir secret sharing over $\mathrm{GF}(257)$}
\label{sec:prelim:shamir}

We instantiate Shamir secret sharing \cite{shamir1979} byte-wise over
$\mathrm{GF}(p)$ with $p = 257$. The choice of $p = 257$ is dictated by
two constraints:
\begin{enumerate}[leftmargin=2em,nosep]
  \item $p > 255$, so every byte value $0,\ldots,255$ is a distinct
        field element.
  \item $p < 2^9$, so each share value fits in a $9$-bit lane, packed
        as one $\mathtt{uint16}$ per byte position on the wire.
\end{enumerate}
The trade-off: $p = 257$ admits at most $p - 1 = 256$ distinct non-zero
evaluation points, capping the committee size at $256$ (constant
\texttt{MaxCommittee257} in \texttt{params.go}). For the Lux validator
committees this is comfortably above the production thresholds (Quasar
operates at $K \le 21$ \cite{lp020}).

\paragraph{Sharing.} To share a $\textsc{seed\_size}$-byte secret
$\mathbf{s} \in [0,256]^{\textsc{seed\_size}}$ across $n$ parties at
threshold $t$:
\begin{enumerate}[leftmargin=2em,nosep]
  \item For each byte position $b \in [0, \textsc{seed\_size})$, sample
        $t-1$ random coefficients
        $a_{b,1}, \ldots, a_{b,t-1} \in \mathrm{GF}(257)$
        and define
        $f_b(x) = \mathbf{s}[b] + \sum_{j=1}^{t-1} a_{b,j} x^j
          \pmod{257}$.
  \item For each recipient $P_i$ at evaluation point $x_i = i$ (one-
        indexed), compute $\mathit{share}_i[b] = f_b(x_i) \bmod 257$
        for all $b$.
  \item Wire-form: pack each $\mathit{share}_i$ as a sequence of
        $\textsc{seed\_size}$ big-endian $\mathtt{uint16}$ values
        ($2 \cdot \textsc{seed\_size}$ bytes total).
\end{enumerate}

\paragraph{Reconstruction.} Given any $t$ shares
$\{(x_i, \mathit{share}_i)\}_{i \in Q}$ with $|Q| \ge t$, the
constant term of each $f_b$ is recovered by Lagrange interpolation at
$x = 0$:
\begin{equation}
  \mathbf{s}[b] \;=\; \sum_{i \in Q} \lambda_i^Q(0) \cdot \mathit{share}_i[b]
  \pmod{257},
  \quad
  \lambda_i^Q(0) = \prod_{j \in Q, j \neq i}
    \frac{-x_j}{x_i - x_j} \bmod 257.
  \label{eq:shamir-lagrange}
\end{equation}
The Lagrange coefficients $\lambda_i^Q(0)$ are computed once per quorum
using Fermat's little theorem for modular inversion
($a^{-1} \equiv a^{p-2} \pmod p$); see \texttt{shamir.go} lines 148--164.

\paragraph{Information-theoretic privacy.} For any $|Q| < t$, the joint
distribution of $\{\mathit{share}_i\}_{i \in Q}$ is uniformly distributed
over $[0,257)^{|Q| \cdot \textsc{seed\_size}}$ \emph{independent} of
$\mathbf{s}$. This is the standard Shamir privacy property
\cite{shamir1979,katzlindell}: every $t-1$-subset of shares reveals
zero information about the secret, as a matter of information theory
(not computational complexity). This property is the load-bearing
secrecy guarantee of Magnetar v0.1 \emph{against the committee} ---
a coalition of $t-1$ or fewer committee members has no useful
information about $S$.

\paragraph{The $\mathrm{GF}(257)$/byte-domain mismatch.} The Lagrange
reconstruction \eqref{eq:shamir-lagrange} produces a value in
$[0, 257)$; the SLH-DSA scheme seed is byte-valued in $[0, 256)$. The
single value $256$ has no byte representation. Magnetar handles this
\emph{constructively}: the Shamir reconstruction is fed into a
cSHAKE256 mix
\begin{equation}
  S = \textsc{cSHAKE256}\bigl(
    \mathrm{byteSum}_\mathit{BE} \,\Vert\, \mathit{committee\_root},
    \mathrm{N}=\text{``Magnetar''},
    \mathrm{S}=\text{``MAGNETAR-SEED-SHARE-V1''},
    8 \cdot \textsc{seed\_size}\,\mathrm{bits}\bigr),
  \label{eq:seed-mix}
\end{equation}
which absorbs the $\mathrm{GF}(257)$-valued $\mathrm{byteSum}$ as a
$2 \cdot \textsc{seed\_size}$-byte big-endian buffer and emits a
$\textsc{seed\_size}$-byte $S$. The cSHAKE256 mix domain-separates the
seed from raw Shamir output and gives a constant-length seed regardless
of the value distribution of $\mathrm{byteSum}$.

\subsection{Threshold signature syntax}
\label{sec:prelim:syntax}

We use the standard threshold-signature syntax of Canetti
\cite{canetti2001uc} and Boneh--Shoup \cite{boneh-shoup-textbook} as
specialised to SLH-DSA.

\begin{definition}[Magnetar threshold signature scheme]
\label{def:magnetar-syntax}
$\Pi_{\mathrm{Magnetar}}$ is a tuple of probabilistic polynomial-time
algorithms $(\textsc{DKG}, \textsc{Sign}, \textsc{Combine},
\textsc{Verify})$:
\begin{itemize}[nosep,leftmargin=2em]
  \item $\textsc{DKG}(1^\lambda, n, t, \mathrm{Committee})$: a
        three-round interactive protocol among committee parties.
        Outputs the joint group public key $pk$, per-party key shares
        $\{(\mathit{NodeID}_i, x_i, \mathit{share}_i)\}_{i=1}^n$, and
        a transcript hash $\tau_{\mathrm{DKG}}$.
  \item $\textsc{Sign}_i(\mathit{share}_i, \mathit{sid}, \kappa,
        \mathrm{Quorum}, \mu)$: an interactive 2-round subprotocol
        executed by party $P_i$ in $\mathrm{Quorum}$. Outputs a
        Round-1 commit message and a Round-2 reveal message.
  \item $\textsc{Combine}(pk, \mu, \mathit{ctx},
        \mathit{sid}, \kappa, \mathrm{Quorum}, t,
        \{\textsc{R1}_i, \textsc{R2}_i\}_{i \in \mathrm{Quorum}},
        \{\mathit{share}_j\}_{j=1}^n)$: a deterministic function
        executed by the aggregator. Outputs a FIPS 205 SLH-DSA
        signature $\sigma$ or $\bot$.
  \item $\textsc{Verify}(pk, \mu, \sigma, \mathit{ctx})$: \emph{the
        unmodified FIPS 205 verifier of \cite{fips205}~\S9.3}.
\end{itemize}
\end{definition}

The deliberate choice that $\textsc{Verify}$ is the unmodified FIPS 205
verifier --- not a Magnetar-specific envelope --- is the load-bearing
operational property: deployment shipping a Magnetar threshold signer
needs to add no new verification code on the consuming side. Every
SLH-DSA verifier in the wild (the FIPS 205 reference, openssl,
cloudflare/circl, libsignal) accepts Magnetar signatures with no
change.

\subsection{Security games}
\label{sec:prelim:games}

We adopt the standard threshold unforgeability game (cf. Pulsar
\cite{lp073pulsarpaper}~\S5).

\begin{definition}[Threshold EUF-CMA]
\label{def:teuf-cma}
$\mathsf{Game}^{\mathrm{tEUF-CMA}}_{\Pi}(\mathcal{A}, \lambda)$:
\begin{enumerate}[nosep,leftmargin=2em]
  \item \textbf{Setup.} The challenger runs $\textsc{DKG}$ with the
        adversary $\mathcal{A}$ playing the corrupt parties.
        $\mathcal{A}$ commits to a corrupt set $F$ with $|F| < t$ at
        the outset (static corruption model). The challenger gives
        $\mathcal{A}$ the joint $pk$ and $\{\mathit{share}_i : i \in F\}$.
  \item \textbf{Queries.} $\mathcal{A}$ adaptively requests signatures
        on messages $\mu_1, \ldots, \mu_{Q_S}$. The challenger runs
        the honest threshold sign protocol and returns each
        $\sigma_i$. $\mathcal{A}$ may invoke the random oracle
        $\mathcal{O}_H$ up to $Q_H$ times.
  \item \textbf{Forgery.} $\mathcal{A}$ outputs $(\mu^\star,
        \sigma^\star)$ with $\mu^\star \notin \{\mu_i\}$, and
        \emph{wins} iff
        $\textsc{Verify}(pk, \mu^\star, \sigma^\star, \mathit{ctx}) = 1$.
\end{enumerate}
The advantage is
$\mathsf{Adv}^{\mathrm{tEUF-CMA}}_{\Pi, \mathcal{A}}(\lambda)
= \Pr[\mathcal{A} \text{ wins}]$.
$\Pi$ is $(Q_S, Q_H, T, \epsilon)$-tEUF-CMA secure if for every
$\mathcal{A}$ within $(Q_S, Q_H, T)$,
$\mathsf{Adv} \le \epsilon$.
\end{definition}

The \emph{aggregator-honest} variant of this game (Magnetar's v0.1
security target) restricts $\mathcal{A}$ from compromising the
aggregator: the aggregator is assumed to execute $\textsc{Combine}$
honestly and zeroize secret material as specified.

\begin{definition}[Aggregator-honest tEUF-CMA]
\label{def:agg-honest-teuf-cma}
The aggregator-honest variant $\mathsf{Game}^{\mathrm{tEUF-CMA-AH}}_{\Pi}
(\mathcal{A}, \lambda)$ is identical to Definition~\ref{def:teuf-cma}
except that the corrupt set $F$ is constrained to exclude the
aggregator process and the adversary has no access to the master seed
$S$ or the reconstructed $sk$ during their in-memory lifetime.
\end{definition}

The v0.1 security claim (Theorem~\ref{thm:magnetar-aggregator-honest}
below) is stated in the aggregator-honest model. The honest disclosure
is that any deployment must operationally enforce
\emph{aggregator-honest} via the mitigations of
\S\ref{sec:intro:trust} or accept the residual risk of an aggregator
compromise. v0.2 (\S\ref{sec:fullmpc}) drops this assumption.
